\documentclass[12font]{article}
\title{\textbf{ Projective Symmetry, Symmetry Anomaly, and Entanglement }}
\author{Chen-Chih Wang  }
\date{\textit{Department of physics, National Tsing Hua University, Hsinchu 30013, Taiwan} \\ (Dated: \today)}
\usepackage[margin=2cm]{geometry}
%\usepackage[fleqn]{amsmath}
%\usepackage{CJK,amsmath}
\usepackage{amsfonts,amssymb}
\usepackage{fancyhdr}
%\pagestyle{fancy}
\usepackage{textcomp}
\usepackage{graphicx}
\usepackage{float}
\usepackage{color}
\usepackage{tikz}
\usepackage{multicol}
\usepackage{threeparttable}
\usepackage{amsthm}
\usepackage{mathtools}
\usepackage{hyperref}
\urlstyle{same}
\usepackage{caption}
\usepackage{subcaption}
\usepackage{dsfont}

\newcommand{\bra}[1]{\left\langle #1 \right|}
\newcommand{\ket}[1]{\left| #1 \right\rangle}
\newcommand{\anglelr}[1]{\left\langle #1 \right\rangle}
\newcommand{\braket}[2]{\left\langle #1 \right|\left. #2 \right\rangle}
\newcommand{\dif}[1]{\frac{\mathrm{d}}{\mathrm{d}#1}}
\newcommand{\diff}[2]{\frac{\mathrm{d}#1}{\mathrm{d}#2}}
\newcommand{\gd}{\boldsymbol{\nabla}}
\newcommand{\dive}{\boldsymbol{\nabla}\cdot}
\newcommand{\curl}{\boldsymbol{\nabla}\times}
\newcommand{\lr}[1]{\left(  #1  \right)}
\newcommand{\lrm}[1]{\left[  #1  \right]}
\newcommand{\lrg}[1]{\left\{  #1  \right\}}
\newcommand{\abs}[1]{\left|  #1  \right|}
\newcommand{\de}{\partial}
\newcommand{\pdif}[1]{\frac{\partial}{\partial#1}}
\newcommand{\pdiff}[2]{\frac{\partial#1}{\partial#2}}
\newcommand{\mrm}[1]{\mathrm{#1}}
\newcommand{\bo}[1]{\mathbf{#1}}
\newcommand{\bog}[1]{\boldsymbol{#1}}
\newcommand{\vint}[3]{\int_{#1}#2\cdot\mathrm{d}\mathbf{#3}}
\newcommand{\voint}[3]{\oint_{#1}#2\cdot\mathrm{d}\mathbf{#3}}
\newcommand{\Vint}{\int\mrm{d}^3\bo{r}'}
\newcommand{\rr}{\mathbf{r} - \mathbf{r}'}
\newcommand{\arr}{\left|\mathbf{r} - \mathbf{r}'\right|}
\newcommand{\alert}[1]{{\color{red}#1}}

\begin{document}
\maketitle

%\fontsize{12pt}{18pt}\selectfont


\begin{abstract}
  
\end{abstract}

\section{Introduction}
The projective symmetry can be used to classify the quantum phase (keywords: 2nd cohomology group, and Ref.\cite{MPS_gapped_phase}). On the other hand, Ref.\cite{Symmetry_and_Entanglement} proved that the state $\ket{\psi}$ in a lattice spin system is long-ranged entangled if 
\[\mathcal{T}\ket{\psi} = e^{iP}\ket{\psi}\quad,\quad\text{where }\;P\neq 0\]
$\mathcal{T}$ is the lattice translation operator. And which can be considered as a projective symmetry 

\alert{Try to relate the \emph{symmetry anomaly} and the \emph{projective symmetry}, they are similar to some extend.}



\begin{thebibliography}{99}

    \bibitem{Symmetry_and_Entanglement}\href{https://doi.org/10.1103/PhysRevX.12.031007}{Lei Gioia and Chong Wang, \emph{Phys. Rev. X} \textbf{12}, 031007 (2022)}
    \bibitem{MPS_gapped_phase}\href{https://doi.org/10.1103/PhysRevB.83.035107}{Xie Chen, Zheng-Cheng Gu, and Xiao-Gang Wen, \emph{Phys. Rev. B} \textbf{83}, 035107 (2011)}

\end{thebibliography}

\end{document}

